\chapter{The Thread Itself}

I said I went to find the Magpie in the morning. That was true, but it was the second true thing. The first was the thread.

It had not gone anywhere in the night. That is the property of thread which makes it better company than people: it waits. The one strand nobody cut still ran out from under the window, silver going grey where it crossed the sill's shadow, and out into a city that walked through it all day without feeling so much as a cobweb. I sat under the window and recalled her teaching: \emph{magic is needlework, Ash. The Weft is every thread of everything, crossing. To follow a thread you do not pull it --- pulling tells the far end. You lay a paw on it, lightly, and you let it tell you which way it is cold. And learn the difference between fraying and cutting, because the world will not teach it to you. Anything frays a thread. Weather frays, and wear frays, and teeth fray, and time frays everything it is left alone with. But cutting one clean is not a skill. It is a right. The Weft knows who held the right, and it keeps the entry either way.}

I put a paw on the line. It was cold toward the city. I went the way it was cold.

The strand ran the rooftops, and it ran them the clean way: chimney to chimney, along ridgelines, over the one plank bridge between the tannery and the ropewalk. Follow a thing's route and you learn the thing. Whoever owned the far end of this thread knew Hollowmere the way I know it: not the streets, which anyone can learn, but the \emph{overstreets}, the roof-roads. I made a note in the ledger, and the note read: \emph{it travels like a cat}, and I did not like the note, so I turned the page.

Halfway across the Shambles' roofline I found the thread being robbed.

A gutter-wisp --- fat, this one, prosperous --- had wound itself around the strand where it sagged between two chimneys, and it was \emph{drinking}: sipping the faint hum off the silver. Not trespass. Grave-robbing. That hum was the last of her moving in the world, and I am not able to write down what I did to that wisp. It went back down its drain unlit. The hum came back to the thread, faint, like a tune through a wall.

The thread dropped to street level at the roofline's end, ran along Needle Lane's gutter --- past the wrong lamps, past Number Twelve, where the dog watched me pass in a silence we had both agreed never to discuss --- and then it turned, and ran straight for the Mere.

And the fog stood up to meet it.

The Mere fog is not the river fog. River fog is weather; it has somewhere else to be. The Mere fog comes up the streets at its own pace, and it does not disperse so much as \emph{decline to leave}, and the lane it occupied that morning was occupied wall to wall, a grey so thick the buildings on either side became opinions. Every window in it dark. Every thread in it slack. The strand ran into that fog low and taut.

There was no fighting through it, because there was nothing to fight, only losing the line to be done. So I did the crossing as she taught me the city must sometimes be done: by paying. A cat's natures are a purse, if you keep them right. Patience, the green, bought me the waiting at each curb while the fog breathed, and near the middle of the lane I felt, for the first time in my life, the bottom of that purse, not empty, but findable. The black bought me the stretches where being seen felt suddenly like a bad idea, though I saw no one.

And twice, at crossings where the cold came up through my feet with that pond-bottom patience I had smelled in her parlor, I spent a little of the pale, her Moonlight, and the cold drew back from me like a tide deciding I was not yet due.

The green comes back with meat and sleep. The pale does not. She stitched it into me, and she is dead, and no cat makes his own --- so every spend from here comes off a bequest, and a bequest runs one way only.

I came out the far side of the fog with wet whiskers, an intact line, and the firm intention of never telling anyone how fast my heart had been going.

The thread ran three more streets, down the terraces, and out onto the sea-wall at the Mereside edge.

And stopped.

Mid-air. Over the water. A cat's height above the black glass of the Mere, the strand ran out past the wall, thinned to a hair, and ended, not frayed, not worn, not melted in the manner of workings that fail. \emph{Cut.} One clean bite, on a slant, the way scissors leave silk. I stood on the sea-wall with my paw on the last of her and looked at that ending a long time.

Understand what it meant. Almost nothing can see a thread. Of the things that can, almost none can touch one, and of those, almost none can cut one clean; her own wards took her forty years and a bone needle to tie and untie. Whoever went out over that water had seen the strand following them, had known it for what it was, and had turned and cut it \emph{one clean stroke}. And they had done it out over open water, where nothing walks.

But it was not only that almost nothing \emph{can} cut a thread. Cutting clean is not a skill; it is a right. The right to cut her work was hers, and she is in the sea. Whoever bit through this strand had no standing to, and the Weft does not punish that and does not forgive it. It only keeps the entry.

And I understood what the strand had been all along. Not her tether to me, that was cut in the parlor, and hangs at my chest still. This was a thread of the \emph{work}: one strand of the forty years the shape had pulled off her walls and wound around itself, snagged on the sill as it left, paying out behind it like a hem coming undone as it walked.

I had not been following her. There is no her to follow. I had been following what \emph{wore} her, and it had felt the line come taut at last, and turned, and bitten through my last road to it.

That is why the strand ran toward the Mere. Everything that night ran toward the Mere. I did not know that yet either.

The Mere lay flat as a ledger to the horizon-fog. Black water; the drowned thing under it that day-side children are told is a story and night-side children are told to stop asking about, sewn shut below the surface by a working I would learn the price of far too late. The fog on the far shore held its shapes as it always does, and one of the shapes, for a moment, was the size of a man, and then was not, and I am a professional, and I wrote it down anyway.

So the thread was not a road. Very well. It was still an answer, and it was this: the killer reads the Weft, travels the overstreets, cuts clean, and did not want to be followed across the water. A trail closed on purpose is a trail worth money, and things worth money have a fence, and the Shambles keeps the best fence in the city.

Which brings me, second, to the Magpie.

\scenebreak

The trouble with visiting the Magpie Exchange is that you must bring something, and the trouble with bringing something, if you are a cat, is that you have no pockets, or rather, you have nine, but I did not intend to spend the contents of my shadow on an introduction. So I did the thing I had been circling since dawn and not doing: I went to her workbox, and I took her second-best thimble.

The second-best. The best is silver, and it stays where her fingers left it, and that is a rule of the new arrangement of the world, and the rule holds. The second-best is brass, worn through at the crown from four decades of the needle's little kiss, worth nothing to anyone who did not watch it work. I carried it to the Shambles in my own mouth, and the whole way there it clicked gently against my teeth like a small tame ghost.

The Shambles after hours is my favorite kind of city: canvas and rope-creak and the day's fish-smell settling like sediment, lanes just wide enough for a barrow and exactly wide enough for me. The Exchange keeps no sign. It keeps a bird.

Brindle was on her perch above the stall shutters, doing the thing magpies do that looks like idleness and is inventory. Black and white with that oil-slick green in the wing. I had seen her at the wake, on the bollard, with her respectful button. She watched me come the length of the lane with the thimble in my teeth and said, before I had set it down:

``You are the witch's cat.''

``Was,'' I said.

She did not say sorry. She turned her head the other way round to look at me again, which in a bird is consideration, and she said, ``Hm,'' in the tone of an appraiser finding the hallmark genuine, and I liked her enormously and permanently.

Business, then. I set the thimble on her counting-board and told her what I wanted to buy: last night but one, Needle Lane, who or what went along it, and anything the lane's gutters had yielded since. Information is the Exchange's true stock; the buttons are a courtesy the trade extends to itself. Brindle looked at the thimble a long time. She did not touch it.

``Payment comes in kinds,'' she said. ``That one is the wrong kind. That one is grief, pet, and I do not fence grief; it only ever comes back on the fencer.'' She hopped once along the rail. ``But work, now. Work I will take. My back alley has grown three nights of wisps --- something has them fat and bold this season, and my stock glimmers, and wisps drink glimmer when the lies run short. Clear my alley and we will talk like professionals.''

So I worked. It took the fat end of the night. Her alley had five of them, coiled in the shutter-slats and the drain-mouths like lit worms in an apple, and the work was low and quiet and mostly patience: a guard here, a cuff there, a cat plainly willing to spend the whole night being unprofitable to drink near. The last and fattest came out of the downspout at the greyest hour and made its argument, which had lights in it, and I declined the argument, and the alley was quiet.

Brindle watched the whole thing from her rail, upside down part of the time, which I chose not to notice.

``You're not useless,'' she said, at dawn. ``That's the whole interview.''

Then she told me what I had bought. A cart on Needle Lane, the night of it: late, unlit, wheels wrapped, and \emph{wax on the wind of it}, she said, wax and something under the wax, pond-water and old rope. Twice along the lane, her pickers had it: up from the city riding light, and down toward the Mere riding heavy, the wrapped wheels telling the stones what the driver would not. I could make nothing of it, and filed it anyway. No pigeons to raise the roosts about it, of course. Of course.

And the lane's gutters had yielded to her pickers, two mornings since, exactly one thing worth a shelf: a scrap of paper, cart-docket shaped, that a gull had beaten her to and a boy had beaten the gull to, and the boy sold to the Exchange for a licorice knot. \emph{Consigned, Needle Lane, night of}, and the consignee line scratched out. Scratched, mind. Not blank. Somebody drove a cart down an unlit street and then did not want to have been the sort of person who had.

She took the thimble at last. ``For holding, pet, not for keeping. Redeem it any dawn.'' She held it a moment longer than trade requires, and turned her head away, and that is as far as a magpie's condolences go, and further than most people's.

I went home along the roofline with a docket-scrap folded into my shadow and the following entry in the ledger: \emph{the night had a cart, the cart had wax, the wax has a guild, and somebody is ashamed of the delivery.} On the chimney-breast board, the candle stub gained a neighbor.

\scenebreak

I lost. In her parlor, on the night, against the thing that wore the room, I lost. It came into her house and it did not even \emph{hurry}. And every plan I had made since --- the thread, the fence, the docket --- ran on ahead of me to the same place: sooner or later I would stand in front of that shape again, and on the current arithmetic it would take everything I am, a second time, at its leisure, while it worked.

So I decided to get better. Deliberately. Like a person. Instead of waiting to be improved by events.

I trained against the goose.

The Watch goose lives behind the tannery and guards six gardens' worth of back-run, and it has never lost an argument in its life, which is precisely what made it worth arguing with. Her rule, from the windowsill years: \emph{the difference between brave and stupid is about four seconds of looking.} So I looked.

I watched what the goose \emph{did} --- not the hissing, which is theater, but the feet, the weight, the neck's coiled bookkeeping --- and I went in low and I lost, and looked, and lost, and looked, and on the fourth argument I did not lose, and left the garden unhissed-at, walking, which no cat in living memory has done.

I trained against the dog at Number Twelve: eleven feet of chain and a standing offer to discuss territory, an offer I had declined daily for years. I accepted the meeting. Slower, that one; the lesson a dog teaches is not speed but \emph{reading}: watch what it will do before deciding what to do about it.

We concluded three sessions. At the end of the third the dog sat down, which for that dog is a treaty, and I sat down too, out of range, as the treaty specified, and we regarded one another with the respect of parties who have wasted none of each other's time. The silence between us on the murder night went undiscussed. Some accounts stay open.

And then the last thing, which was not fighting, and which was the reason for all the fighting.

Her ward is still on the parlor door frame. It hangs there with a hole in it the size of what walked through, the stitches lifted, not torn, a wound with straight edges. I had watched her patch that ward four hundred times: after storms, after solstices, once after an argument with the chimney that I maintain to this day the chimney started. Watching is not doing. I know that now better than any creature alive.

I got up on the frame with a mouthful of her thread from the workbox, and I mended it.

I will not pretend it was her work. My stitches sat in the old ward like a boy's writing on a master's page --- fat, anxious, twice-knotted where hers would have held with a breath. It took most of a night, and one stitch I placed and unpicked eleven times because it \emph{pulled the wrong way}; my paws wanted to run the line mirror-wise, right to left, and the ward wanted otherwise, and I remember being irritated at my own hands and thinking nothing of it. Nothing in this book costs me more.

But when it was finished the ward held. The door frame hummed again, low, the sound of a house with someone minding it, and the hole the size of what walked through was gone.

Her work always took its payment at the pull. I had watched that for fifteen years without giving it a name: she set a stitch, and something went out of her with the thread. A warmth. A face she used to be able to call up. An hour of luck she had been keeping. A year off the far end.

Nobody had told me what a cat pays with. I found out at about the fourth stitch. A familiar has one coin, and the coin is sleep. I paid a night of it into that frame, and the frame was not done with me, so I paid the next night as well: two of them on the cold hearthstone with my eyes open on nothing while the dark went about its business without me.

I sat on the hearthstone and looked at my crooked stitches in her doorway and understood that I was going to do all of it, the whole case, every thread of it, and that no one was going to stop me, including whatever had looked at me and decided I was not an interruption.

It would learn. I was exactly an interruption.

\scenebreak

That night I walked the rounds nobody was paying for: the lamps still short-handed, past the drains and the dark doorways, down Needle Lane to its cold end and back. Two wisps put up their small arguments and lost them. The dog at Number Twelve stood, once, and did not bark --- looked where I looked, at the fog standing at the lane's end where fog does not stand --- and we held that watch together for a while, chain and cat, in the silence of colleagues.

Wisps used to be rare on this street. Used to be. The dark had noticed the missing lamps, the way water notices a low place. Somebody had to walk the line while the Lamplighters grieved.

The city I can work in is a particular kind of city. I intended to keep it.
